\documentclass[a4paper,12pt]{article}
\usepackage[english,russian]{babel}
\usepackage[utf8]{inputenc}
\usepackage{latexsym,mathrsfs}
\usepackage{stmaryrd, enumitem}
\usepackage{amsthm,amsfonts,amssymb,amsmath}
\usepackage{geometry}
\usepackage{tempora}
\usepackage{graphicx}
\usepackage{url}
\usepackage{listingsutf8}


\geometry{top=17mm}  \geometry{bottom=20mm}
\geometry{left=17mm} \geometry{right=17mm}
\linespread{1.1}
\parindent=5mm

\def\titleK#1{\begin{center}{\Large\textbf{#1}}\end{center}}
\def\authorK#1{\begin{center}{#1}\end{center}}

\newenvironment{abstractK}{}

\newtheorem{theorem}{Теорема}
\newtheorem{lemma}{Лемма}
\newtheorem{corollary}{Следствие}
\newtheorem{proposition}{Предложение}
\newtheorem{statement}{Утверждение}
\theoremstyle{definition}
\newtheorem{definition}{Определение}
\newtheorem{problem}{Задача}
\newtheorem{question}{Вопрос}
\newtheorem{conjecture}{Гипотеза}
\newtheorem{construction}{Конструкция}
\newtheorem{remark}{Замечание}

%ВАЖНО: Просьба пользоваться имеющимися окражениями-теоремами
%ВАЖНО: Не менять и не добавлять ничего выше этой строки.
%Изменения вносить только внутри окружения \begin{document}\end{document}

\begin{document}

%Название доклада
\titleK{Исследование аппаратно-программных реализаций алгоритмов защиты информации}

%Информация об авторах (ЗАМЕНИТЕ на свои данные)
\authorK{Д.~С.~Соловьёв$^{1}$, А. Р. Труфанов$^{2}$, С. В. Копылова$^{2}$, Д. А. Бондаренко$^{2}$,  А. С. Марков$^{2}$, 
\\В. М. Завьялова$^{3}$, Э. А. Салеев$^{4}$, А. В. Тюндин$^{5}$
\\$^{1}$Южный федеральный университет 
\\$^{2}$НИЯУ МИФИ
\\$^{3}$СПбГУ
\\$^{4}$Казанский федеральный университет
\\$^{5}$НГТУ им. Р. Е. Алексеева
\\{\bf E-mail:} dsolove@sfedu.ru, 111@mail.ru,  SonyCop26@yandex.ru, 111@mail.ru, 111@mail.ru, 111@mail.ru, 111@mail.ru, AlexanderTyundin@yandex.ru,}

\begin{abstract}
      Аннотация. (Ждем ответ от Яковлева А.А.)\\
\\{\bf Ключевые слова:} \textit{реверс-инжиниринг, статический и динамический анализ, фреймворк Qt, поточное шифрование, антиотладка, обход защиты, информационная безопасность. (пересмотреть ключевые слова)} 
\end{abstract}

\begin{abstractK}


\section*{Введение}

В современных программных продуктах, предназначенных для защиты коммерческой информации и интеллектуальной собственности, широко применяются механизмы противодействия исследованию исполняемого кода. Наиболее распространёнными из них являются динамическая обфускация, аппаратная защита криптографических операций и методы противодействия статическому и динамическому анализу.

Объектом исследования является исполняемое приложение, использующее комплекс механизмов защиты программного кода. Цель работы заключается в выявлении применяемых механизмов защиты, исследовании используемого обратимого криптографического преобразования и формировании рекомендаций по анализу аналогичных программных систем.

Для достижения поставленной цели использованы методы статического и динамического анализа исполняемого кода.

\section*{1. Методика исследования}

Исследование выполнялось в два этапа.

На первом этапе средствами статического анализа определялись применяемые механизмы защиты исполняемого кода, структура взаимодействия программных модулей и характер выполняемых преобразований.

На втором этапе выполнялось динамическое исследование процесса исполнения программы, позволившее определить последовательность восстановления защищённых участков кода, выявить механизмы противодействия отладке и установить назначение отдельных функциональных модулей.

\section*{2. Результаты исследования механизмов защиты}

\subsection*{2.1. Cryptographic Code Protection}

Средствами статического и динамического анализа установлено, что исследуемое приложение использует технологию \textit{Cryptographic Code Protection} (динамическую обфускацию исполняемого кода), реализованную программным протектором.

Применяемый механизм основан на частичном сокрытии исполняемого кода. До начала выполнения защищённые участки отсутствуют в открытом виде и восстанавливаются непосредственно в адресном пространстве процесса с использованием аппаратного средства защиты информации. Для локализации защищённых областей используются специальные сигнатуры, после обнаружения которых выполняется обратимое криптографическое преобразование соответствующего участка памяти.

Исследование показало, что после восстановления исполняемого кода в него дополнительно интегрируются механизмы контроля среды выполнения. К ним относятся проверки наличия внешнего отладчика, контроль целостности исполняемого образа, обнаружение признаков выполнения программы в виртуализированной среде, а также механизмы, допускающие виртуализацию отдельных функций. Таким образом, защита реализуется не только посредством сокрытия исполняемого кода, но и за счёт встроенных средств противодействия исследованию.

Использование динамической обфускации существенно осложняет проведение статического анализа, поскольку значительная часть исполняемого кода становится доступной только во время выполнения программы.

\subsection*{2.2. Противодействие динамическому анализу}

В результате исследования установлено, что приложение использует несколько взаимодополняющих механизмов противодействия динамическому анализу.

Одним из них является технология \textit{Trap-Based}, основанная на выполнении привилегированных операций с последующим анализом способа обработки возникающих исключений. Такой подход позволяет выявлять вмешательство отладчика или изменение стандартного механизма обработки исключений.

Дополнительно выявлено применение технологии \textit{Self-Debugging}, основанной на монополизации порта отладки собственным процессом, что затрудняет подключение внешнего отладчика.

Кроме того, средствами статического анализа установлено наличие дополнительных проверок среды выполнения. Обнаружены обращения к флагам \texttt{BeingDebugged} и \texttt{NtGlobalFlag} структуры \texttt{PEB}, а также косвенный вызов функции проверки через виртуальную таблицу класса, что затрудняет её обнаружение при анализе графа вызовов.

Совместное применение технологий \textit{Trap-Based}, \textit{Self-Debugging} и механизмов динамической обфускации существенно повышает устойчивость приложения к проведению динамического анализа.

\section*{3. Исследование криптографической подсистемы}

Анализ восстановленного исполняемого кода показал наличие отдельной подсистемы, выполняющей обратимое криптографическое преобразование данных.

Установлено, что преобразование осуществляется с использованием аппаратного средства защиты информации посредством специализированного программного интерфейса. Перед выполнением операции приложение проходит процедуру аутентификации на аппаратном устройстве, после чего выполняется обработка данных с использованием закрытого криптографического материала, недоступного процессу в открытом виде.

Исследование внутренней реализации показало наличие двух массивов внутреннего состояния объёмом по 512 элементов, которые используются при генерации выходной последовательности. Обновление внутреннего состояния сопровождается операциями циклического сдвига, сложения по модулю $2^{32}$ и исключающего ИЛИ, что характерно для потоковых криптографических алгоритмов.

По совокупности выявленных признаков можно предположить, что исследуемое обратимое преобразование построено на криптографическом примитиве, архитектурно близком к семейству HC-256. Данный вывод носит вероятностный характер, поскольку часть внутренней логики выполняется аппаратным средством защиты информации и недоступна для непосредственного анализа.

Для идентификации используемого преобразования применялась совокупность структурных признаков: объём внутреннего состояния, организация внутренних буферов, используемые побитовые операции и характер обновления внутреннего состояния. Такой подход позволяет выполнять предварительную идентификацию криптографических примитивов даже при отсутствии полной информации об их реализации.

\section*{4. Практические рекомендации}

Проведённое исследование показало, что выполнение динамического анализа подобных приложений возможно после локализации механизмов противодействия отладке и последующего исследования уже восстановленного исполняемого кода.

Для анализа программ, использующих аналогичные методы защиты, рекомендуется применять поэтапный подход, включающий выявление механизмов динамической обфускации, исследование взаимодействия с аппаратными средствами защиты информации и анализ восстановленного исполняемого кода.

Предложенный подход позволяет существенно сократить объём исследуемого кода и сосредоточить анализ на участках программы, содержащих основную функциональную нагрузку.

\section*{Заключение}

В ходе исследования были выявлены основные механизмы защиты исследуемого приложения, выполнен анализ динамической обфускации исполняемого кода, исследован механизм использования аппаратного средства защиты информации при выполнении обратимого криптографического преобразования данных и проведена предварительная идентификация применяемого криптографического примитива.

Полученные результаты могут быть использованы при исследовании программных продуктов, использующих аналогичные механизмы защиты исполняемого кода, а также при разработке методик анализа защищённых программных систем.
\end{document}
